%% sec4_algorithms.tex -- Section 4: Algorithms

\section{Algorithms}
\label{sec:algorithms}

We describe two procedures for constructing a multi-ticket collection with
small total pairwise overlap.  The first builds the collection greedily by
maximising coverage; the second refines the greedy output via simulated
annealing.

\subsection{Greedy Coverage Selection}
\label{sec:algo-greedy}

The maximum coverage problem asks: given a universe $U$, a family
$\mathcal{F}$ of subsets of $U$, and a budget $n$, choose $n$ sets from
$\mathcal{F}$ to maximise $|\bigcup_{i=1}^{n} A_i|$.  \citet{Nemhauser1978}
proved that the greedy algorithm, which at each step adds the set with the
largest marginal coverage gain, achieves a $(1-1/e)$-approximation.
\citet{Feige1998} showed this ratio is tight under standard complexity
assumptions.

Algorithm~\ref{alg:greedy} instantiates this paradigm with $U = [v]$ and
$\mathcal{F} = \binom{[v]}{k}$.  Maximising coverage of $[v]$ by the selected
tickets minimises the number of elements counted more than once, which
corresponds to minimising $\Omega(\mathcal{A})$.  A tie-breaking rule based
on cumulative overlap makes the secondary objective explicit.

\begin{algorithm}[H]
  \caption{Greedy Coverage Selection}
  \label{alg:greedy}
  \KwInput{Universe size $v$, ticket size $k$, number of tickets $n$,
           candidate pool $\mathcal{P} \subseteq \binom{[v]}{k}$
           of size $M \gg n$.}
  \KwOutput{Collection $\mathcal{A} = (A_1, \ldots, A_n)$ with small
            total overlap $\Omega(\mathcal{A})$.}
  \BlankLine
  $\mathrm{covered} \leftarrow \emptyset$\;
  $\mathrm{cumOv}[A] \leftarrow 0$ for all $A \in \mathcal{P}$\;
  $\mathcal{A} \leftarrow ()$\;
  \BlankLine
  \For{$\ell = 1$ \KwTo $n$}{
    $\mathrm{gain}[A] \leftarrow |A \setminus \mathrm{covered}|$
    for all available $A \in \mathcal{P}$\;
    \BlankLine
    $\mathrm{score}[A] \leftarrow
      (L+1) \cdot \mathrm{gain}[A] - \mathrm{cumOv}[A]$
    \tcp*[r]{$L = n \cdot k$: ensures gain dominates overlap}
    \BlankLine
    $A^* \leftarrow \arg\max_{A \in \mathcal{P} \setminus \mathcal{A}}
      \mathrm{score}[A]$\;
    \BlankLine
    Append $A^*$ to $\mathcal{A}$\;
    $\mathrm{covered} \leftarrow \mathrm{covered} \cup A^*$\;
    $\mathrm{cumOv}[A] \leftarrow \mathrm{cumOv}[A] + |A \cap A^*|$
    for all $A \in \mathcal{P}$\;
  }
  \Return $\mathcal{A}$\;
\end{algorithm}

\begin{remark}[Scoring rule]
  The multiplier $L+1 = nk+1$ in Algorithm~\ref{alg:greedy} guarantees that
  any non-zero coverage gain makes a ticket strictly preferred over any ticket
  with smaller gain, regardless of overlap.  Formally, if $\mathrm{gain}[A]
  > \mathrm{gain}[B]$, then $\mathrm{score}[A] > \mathrm{score}[B]$ for
  any values of $\mathrm{cumOv}$ bounded above by $Lk$ (since
  $\mathrm{gain}$ takes integer values in $[0, v]$, one unit of gain is
  worth $(L+1) > Lk$ tie-breaker units).  Overlap is therefore a
  \emph{secondary} objective, used only when coverage gains are equal.
\end{remark}

\begin{theorem}[Approximation guarantee]
  \label{thm:greedy-guarantee}
  Let $\mathrm{OPT}_n = \max_{\mathcal{A}: |\mathcal{A}|=n}
  |\bigcup_{A \in \mathcal{A}} A|$ be the maximum coverage achievable with
  $n$ tickets chosen from the full family $\binom{[v]}{k}$.
  Algorithm~\ref{alg:greedy} run on the full family (i.e., $M =
  \binom{v}{k}$) produces $\mathcal{A}$ with
  \[
    \left|\bigcup_{A \in \mathcal{A}} A\right|
    \;\ge\; \left(1 - e^{-1}\right) \cdot \mathrm{OPT}_n
    \;\approx\; 0.632 \cdot \mathrm{OPT}_n.
  \]
\end{theorem}

\begin{proof}
  The family $\binom{[v]}{k}$ constitutes a set system on the ground set
  $[v]$, and Algorithm~\ref{alg:greedy} is exactly the greedy algorithm for
  the maximum-coverage problem on this system (with $|\mathcal{A}| = n$ as
  the cardinality constraint).  The $(1-1/e)$ bound follows from Theorem~2
  of Nemhauser et al.\ \cite{Nemhauser1978}, which applies to any monotone
  submodular function; coverage is the canonical example.
\end{proof}

\begin{remark}[Practical pool size]
  Running Algorithm~\ref{alg:greedy} on the full family $\binom{[v]}{k}$
  is infeasible for large parameters (e.g., $\binom{25}{15} =
  3{,}268{,}760$ for the Lotofácil).  In practice we draw a random pool
  $\mathcal{P}$ of size $M = \max(500, 20n)$ and run the algorithm on
  $\mathcal{P}$.  The approximation guarantee no longer applies rigorously,
  but empirical evidence (Section~\ref{sec:empirical}) shows the resulting
  collections are near-optimal.  For the regime $v=25$, $k=15$, the greedy
  construction typically achieves full coverage of $[25]$ after just
  $\lceil v/k^* \rceil = 2$ tickets (where $k^* = k - m^* = 10$ is the
  number of unique elements per ticket), so coverage saturates rapidly and
  the tie-breaking overlap rule dominates for $n \ge 2$.
\end{remark}

\subsection{Simulated Annealing Refinement}
\label{sec:algo-sa}

The greedy solution is not guaranteed to minimise $\Omega(\mathcal{A})$
globally, since each greedy step is locally optimal.  We apply Simulated
Annealing (SA) \cite{Kirkpatrick1983} to refine the greedy solution by
minimising the total pairwise overlap objective
\[
  \Omega(\mathcal{A}) = \sum_{1 \le i < j \le n} |A_i \cap A_j|.
\]

\begin{algorithm}[H]
  \caption{Simulated Annealing for Overlap Minimisation}
  \label{alg:sa}
  \KwInput{Initial collection $\mathcal{A}^{(0)}$ (greedy output),
           iteration count $N_\text{iter}$,
           initial temperature $T_0$, cooling factor $\alpha \in (0,1)$.}
  \KwOutput{Refined collection $\mathcal{A}^*$ with $\Omega(\mathcal{A}^*)
            \le \Omega(\mathcal{A}^{(0)})$ (with high probability).}
  \BlankLine
  $\mathcal{A} \leftarrow \mathcal{A}^{(0)}$;
  $\mathcal{A}^* \leftarrow \mathcal{A}$;
  $T \leftarrow T_0$\;
  Compute $\bm{b}[A_i] \in \{0,1\}^v$ (indicator vector) for each $A_i$\;
  \BlankLine
  \For{$s = 1$ \KwTo $N_\text{iter}$}{
    Choose $i \in [n]$ and $j \in [k]$ uniformly at random\;
    Choose $x \in [v] \setminus A_i$ uniformly at random\;
    $A_i' \leftarrow (A_i \setminus \{A_i[j]\}) \cup \{x\}$\;
    \BlankLine
    $\delta \leftarrow
      \sum_{\ell \ne i} |A_\ell \cap A_i'|
      - \sum_{\ell \ne i} |A_\ell \cap A_i|$
    \tcp*[r]{Efficient $O(nv)$ delta via indicator vectors}
    \BlankLine
    \If{$\delta < 0$ \textbf{or} $\mathrm{Uniform}(0,1) < e^{-\delta/T}$}{
      $A_i \leftarrow A_i'$;
      update $\bm{b}[A_i]$;
      $\Omega \leftarrow \Omega + \delta$\;
      \If{$\Omega < \Omega(\mathcal{A}^*)$}{
        $\mathcal{A}^* \leftarrow \mathcal{A}$\;
      }
    }
    $T \leftarrow \alpha \cdot T$\;
  }
  \Return $\mathcal{A}^*$\;
\end{algorithm}

\begin{remark}[Efficient delta computation]
  \label{rem:efficient-delta}
  The overlap increment $\delta$ in Algorithm~\ref{alg:sa} is computed in
  $O(nv)$ time using boolean indicator vectors $\bm{b}[A_i] \in \{0,1\}^v$,
  maintained incrementally: $\sum_{\ell \ne i} |A_\ell \cap A_i'| =
  \sum_{\ell \ne i} \langle \bm{b}[A_\ell], \bm{b}[A_i']\rangle$.
  Only the row corresponding to the modified ticket needs updating,
  reducing the per-iteration cost from $O(n^2 v)$ to $O(nv)$.
  Incremental evaluation of this kind is standard practice in local-search
  implementations; see \citet{Talbi2009}, Section~2.1.4.
\end{remark}

\subsection{Complexity analysis}
\label{sec:algo-complexity}

\begin{proposition}[Algorithm complexities]
  \label{prop:complexity}
  With pool size $M = O(nk)$ and $N_\mathrm{iter}$ SA iterations,
  pool generation costs $O(Mv\log v)$, greedy selection costs
  $O(n \cdot M \cdot v)$ \citep{Johnson1974}, and each SA run costs
  $O(N_\mathrm{iter} \cdot n \cdot v)$ (Remark~\ref{rem:efficient-delta}).
  For the Lotof\'{a}cil with $v=25$, $k=15$, $n \le 100$, $M=500$,
  $N_\mathrm{iter}=2000$, the total cost per run is $O(10^6)$ operations,
  completing in under one second on modern hardware.
  The complexity is linear in $n$ and constant in $v$ for fixed $v=25$,
  so the algorithms scale without difficulty to large budgets.
\end{proposition}

\begin{remark}[NP-hardness does not apply here]
  The maximum-coverage problem on a general set system is NP-hard
  \cite{Feige1998}.  For our specific structure, however, the universe has
  fixed size $v=25$, so exhaustive enumeration over $\binom{v}{k}^n$ is
  polynomial in $n$ for fixed $v$.
  The greedy approximation is used for speed, not necessity.  In the
  Lotofácil instance, the greedy solution is very close to optimal:
  full coverage of $[25]$ is achievable with $n = 2$ tickets (since
  $2 \times (k - m^*) = 2 \times 10 = 20 > v - k = 10$), so for $n \ge 2$
  the greedy always achieves $|\bigcup A_i| = v = 25$, and the SA's role
  is purely to minimise $\Omega$ within the set of maximum-coverage solutions.
\end{remark}
